% Options for packages loaded elsewhere
\PassOptionsToPackage{unicode}{hyperref}
\PassOptionsToPackage{hyphens}{url}
\PassOptionsToPackage{table}{xcolor}
%
\documentclass[
  11pt,
]{report}
\usepackage{amsmath,amssymb}
\usepackage{iftex}
\ifPDFTeX
  \usepackage[T1]{fontenc}
  \usepackage[utf8]{inputenc}
  \usepackage{textcomp} % provide euro and other symbols
\else % if luatex or xetex
  % pdfLaTeX-only repository policy: the original unicode-math/fontspec
  % branch is unreachable under pdfLaTeX and is replaced with a guard.
  \PackageError{cpython-suite}{Compile with pdfLaTeX only}{}
\fi
\usepackage{lmodern}
\ifPDFTeX\else
  % xetex/luatex font selection
\fi
% Use upquote if available, for straight quotes in verbatim environments
\IfFileExists{upquote.sty}{\usepackage{upquote}}{}
\IfFileExists{microtype.sty}{% use microtype if available
  \usepackage[]{microtype}
  \UseMicrotypeSet[protrusion]{basicmath} % disable protrusion for tt fonts
}{}
\makeatletter
\@ifundefined{KOMAClassName}{% if non-KOMA class
  \IfFileExists{parskip.sty}{%
    \usepackage{parskip}
  }{% else
    \setlength{\parindent}{0pt}
    \setlength{\parskip}{6pt plus 2pt minus 1pt}}
}{% if KOMA class
  \KOMAoptions{parskip=half}}
\makeatother
\usepackage{xcolor}
\usepackage{listings}
\newcommand{\passthrough}[1]{#1}
\lstset{defaultdialect=[5.3]Lua}
\lstset{defaultdialect=[x86masm]Assembler}
\usepackage{longtable,booktabs,array}
\usepackage{calc} % for calculating minipage widths
% Correct order of tables after \paragraph or \subparagraph
\usepackage{etoolbox}
\makeatletter
\patchcmd\longtable{\par}{\if@noskipsec\mbox{}\fi\par}{}{}
\makeatother
% Allow footnotes in longtable head/foot
\IfFileExists{footnotehyper.sty}{\usepackage{footnotehyper}}{\usepackage{footnote}}
\makesavenoteenv{longtable}
\setlength{\emergencystretch}{3em} % prevent overfull lines
\providecommand{\tightlist}{%
  \setlength{\itemsep}{0pt}\setlength{\parskip}{0pt}}
\setcounter{secnumdepth}{5}
% CPython Engineering Suite shared presentation layer -- pdfLaTeX compatible
\usepackage[T1]{fontenc}
\usepackage{lmodern}
\usepackage{microtype}
\usepackage{geometry}
\geometry{letterpaper,margin=0.91in,headheight=16pt,headsep=19pt}
\usepackage{xcolor} % [table] passed at top of file (option-clash fix)
\definecolor{CPNavy}{HTML}{13283B}
\definecolor{CPBlue}{HTML}{176BA0}
\definecolor{CPGold}{HTML}{D8A73C}
\definecolor{CPMuted}{HTML}{566979}
\definecolor{CPBg}{HTML}{EFF5F8}
\usepackage{titlesec}
\titleformat{\chapter}[display]{\normalfont\sffamily\bfseries\color{CPNavy}}{\large\color{CPBlue}CHAPTER \thechapter}{9pt}{\Huge}[\vspace{6pt}\color{CPGold}\titlerule]
\titleformat{\section}{\Large\bfseries\sffamily\color{CPBlue}}{\thesection}{0.7em}{}
\titleformat{\subsection}{\large\bfseries\sffamily\color{CPNavy}}{\thesubsection}{0.7em}{}
\usepackage{fancyhdr}
\pagestyle{fancy}
\fancyhf{}
\fancyhead[L]{\footnotesize\color{CPMuted}\sffamily CPYTHON ENGINEERING SUITE}
\fancyhead[R]{\footnotesize\color{CPMuted}\sffamily 3.16.0a1+dev SNAPSHOT}
\fancyfoot[C]{\footnotesize\color{CPMuted}\thepage}
\renewcommand{\headrulewidth}{0.3pt}
\AtBeginDocument{\urlstyle{tt}}
\lstset{basicstyle=\ttfamily\small,columns=fullflexible,keepspaces=true,breaklines=true,breakatwhitespace=false,showstringspaces=false,frame=single,framesep=5pt,rulecolor=\color{CPBg},backgroundcolor=\color{CPBg},xleftmargin=3pt,xrightmargin=3pt}
\usepackage{tcolorbox}
\tcbuselibrary{breakable,skins}
\newtcolorbox{sourcebox}[1][]{enhanced,breakable,sharp corners,boxrule=0.4pt,colback=CPBg,colframe=CPBlue,borderline west={3pt}{0pt}{CPBlue},left=10pt,right=8pt,top=6pt,bottom=6pt,#1}
\newtcolorbox{cautionbox}[1][]{enhanced,breakable,sharp corners,boxrule=0.4pt,colback=CPBg,colframe=CPGold,borderline west={3pt}{0pt}{CPGold},left=10pt,right=8pt,top=6pt,bottom=6pt,#1}
\usepackage{tikz}
\usetikzlibrary{arrows.meta,positioning,shapes.geometric,calc}
\usepackage{enumitem}
\setlist{itemsep=2pt,topsep=4pt}
\usepackage{booktabs}
\usepackage{longtable}
\AtBeginEnvironment{longtable}{\small}
\usepackage{array}
\setlength{\parskip}{5pt}
\setlength{\parindent}{0pt}
\setcounter{tocdepth}{2}
\usepackage{hyperref}
\hypersetup{colorlinks=true,linkcolor=CPBlue,urlcolor=CPBlue,citecolor=CPBlue,pdfauthor={CPython Engineering Documentation Suite}}
\makeatletter
\renewcommand{\maketitle}{%
  \begin{titlepage}\pagecolor{CPNavy}\color{white}\sffamily
  \vspace*{0.8in}
  {\color{CPGold}\rule{2.1in}{4pt}}\\[16pt]
  {\Large\bfseries CPYTHON ENGINEERING DOCUMENTATION SUITE}\\[24pt]
  {\Huge\bfseries \@title\par}\vspace{19pt}
  \vfill
  {\normalsize Source baseline: CPython 3.16.0 alpha 1+dev repository snapshot}\par
  {\normalsize Prepared 16 October 2026}\par
  {\small\color{white!70} Source-traced engineering commentary; not an official Python specification.}\vspace{0.65in}
  \end{titlepage}\nopagecolor\color{black}
  }
\makeatother
\ifLuaTeX
  \usepackage{selnolig}  % disable illegal ligatures
\fi
\usepackage{bookmark}
\IfFileExists{xurl.sty}{\usepackage{xurl}}{} % add URL line breaks if available
\urlstyle{same}
\hypersetup{
  pdftitle={CPython Engineering Curriculum and Laboratory},
  hidelinks,
  pdfcreator={LaTeX via pandoc}}

\title{CPython Engineering Curriculum and Laboratory}
\usepackage{etoolbox}
\makeatletter
\providecommand{\subtitle}[1]{% add subtitle to \maketitle
  \apptocmd{\@title}{\par {\large #1 \par}}{}{}
}
\makeatother
\subtitle{Volume 08 \textbar{} A prerequisite-ordered twelve-module
study and implementation program}
\author{}
\date{}

\begin{document}
\maketitle

{
\setcounter{tocdepth}{1}
\tableofcontents
}
\chapter{Program charter and learner
outcomes}\label{program-charter-and-learner-outcomes}

\begin{sourcebox}
\textbf{Evidence classification.} \textbf{SOURCE} = statements supported by material present in the uploaded snapshot; \textbf{IMPLEMENTATION} = CPython-specific internals, not guaranteed Python semantics; \textbf{ENGINEERING} = original recommendations, examples, checks, and exercises. The source baseline is \textbf{3.16.0 alpha 1+dev}; missing repository file bodies are noted as evidence limitations.
\end{sourcebox}

This curriculum integrates Volumes 01--07 into a reproducible technical
training path. Its goals are to understand Python language semantics,
navigate the CPython implementation, trace the compiler and interpreter,
design safe native extensions, analyze concurrency/lifetime behavior,
and validate software with security and performance discipline. Module
sequence and grading rubric are \textbf{ENGINEERING recommendations}.
Repository descriptions and paths come from the uploaded 3.16.0 alpha
1+dev snapshot.

\section{Competency ladder}\label{competency-ladder}

\begin{longtable}[]{@{}
  >{\raggedright\arraybackslash}p{(\columnwidth - 4\tabcolsep) * \real{0.3333}}
  >{\raggedright\arraybackslash}p{(\columnwidth - 4\tabcolsep) * \real{0.3333}}
  >{\raggedright\arraybackslash}p{(\columnwidth - 4\tabcolsep) * \real{0.3333}}@{}}
\toprule\noalign{}
\begin{minipage}[b]{\linewidth}\raggedright
Tier
\end{minipage} & \begin{minipage}[b]{\linewidth}\raggedright
Competency
\end{minipage} & \begin{minipage}[b]{\linewidth}\raggedright
Observable evidence
\end{minipage} \\
\midrule\noalign{}
\endhead
\bottomrule\noalign{}
\endlastfoot
Foundation & Explain Python syntax, binding, objects and exceptions &
Tests and semantic trace \\
Application & Select appropriate library interfaces and validate data &
Bounded, well-tested service \\
Implementation & Explain compiler, code objects, frames and evaluation &
Stage-by-stage annotated trace \\
Native & Build and audit a safe C extension & Ownership/error-path
checklist \\
Systems & Reason about GC, free threads and reclamation & Race/lifetime
investigations \\
Assurance & Create trustworthy build and runtime evidence & CI,
benchmark, threat model \\
\end{longtable}

\chapter{Prerequisite graph and recommended
sequence}\label{prerequisite-graph-and-recommended-sequence}

\begin{lstlisting}
M01 Python execution & syntax
  |\
  | +-- M02 objects, data model and exceptions
  |       +-- M03 library design and trustworthy I/O
  |               +-- M04 testing and security foundations
  |
  +-- M05 CPython source navigation
          +-- M06 tokenizer, PEG and AST
                  +-- M07 compiler, CFG and code objects
                          +-- M08 frames, interpreter, specialization
                                  +-- M09 GC, ownership and threads
                                          +-- M10 C API and extension modules
                                                  +-- M11 diagnostics & performance
                                                        +-- M12 capstone review
\end{lstlisting}

Cross-dependency: M04 must be completed before M10, and M09 before M10;
the program is not a single uninterrupted linear chain. Learners
comfortable with language semantics may challenge M01--M03 by producing
equivalent artifacts, but the C API and memory modules should not be
skipped without demonstrated ownership knowledge.

\chapter{Module specifications}\label{module-specifications}

\section{M01 --- Language execution
model}\label{m01-language-execution-model}

\textbf{Prerequisite:} basic programming experience. \textbf{Reading:}
Volume 01;
\path{Doc/reference/lexical_analysis.rst},
\path{executionmodel.rst};
\path{Doc/tutorial/controlflow.rst}.
\textbf{Objectives:} explain the difference between lexical structure,
binding and execution; predict what a function returns before running
it; distinguish local, nonlocal and global names.

\textbf{Lab:} create a scope exploration suite with closures and a
deliberately unbound local. Write one test per expected behavior and
draw a name-resolution diagram. \textbf{Exit criterion:} at least ten
assertions and an accurate explanation of why assignments rebind names
instead of copying referenced objects.

\section{M02 --- Objects, protocols and
exceptions}\label{m02-objects-protocols-and-exceptions}

\textbf{Prerequisite:} M01. \textbf{Reading:} Volume 01 object/exception
chapters; \path{Doc/tutorial/classes.rst},
\passthrough{\lstinline!errors.rst!}. \textbf{Objectives:} implement
iteration, context management, immutable-style data objects and
exception chaining.

\textbf{Lab:} build a resource-owning iterator that closes its
underlying file on normal and error paths. \textbf{Exit criterion:}
tests demonstrate safe cleanup, correct use of identity vs equality and
an explicit statement that garbage-collector timing is not a
language-level cleanup strategy.

\section{M03 --- Standard-library
engineering}\label{m03-standard-library-engineering}

\textbf{Prerequisite:} M01 and M02. \textbf{Reading:} Volume 02;
\path{Doc/library/json.rst},
\passthrough{\lstinline!asyncio.rst!},
\passthrough{\lstinline!venv.rst!},
\passthrough{\lstinline!secrets.rst!},
\path{security_warnings.rst}. \textbf{Objectives:}
validate decoded data, manage paths, distinguish async/thread/process
designs, keep secrets out of logs.

\textbf{Lab:} implement a bounded JSON ingestion pipeline and write an
operational configuration guide. \textbf{Exit criterion:} valid and
invalid input tests; deterministic output replacement; managed task
cancellation; resource limits stated explicitly.

\section{M04 --- Testing and threat
boundaries}\label{m04-testing-and-threat-boundaries}

\textbf{Prerequisite:} M02 and M03. \textbf{Reading:} Volume 07;
\path{Doc/library/pickle.rst},
\path{security_warnings.rst}; relevant
\passthrough{\lstinline!Lib/test/!} modules. \textbf{Objectives:} apply
unit/integration testing, classify trust boundaries, design negative
tests, and perform release readiness reviews.

\textbf{Lab:} add malformed-input tests, timeout tests and an
adversarial parameter test to the ingestion pipeline. \textbf{Exit
criterion:} no untrusted pickle execution, no shell command string
concatenation and all failing tests block packaging.

\section{M05 --- Repository navigation and
architecture}\label{m05-repository-navigation-and-architecture}

\textbf{Prerequisite:} M01. \textbf{Reading:} Volume 03;
\path{InternalDocs/README.md},
\passthrough{\lstinline!structure.md!},
\passthrough{\lstinline!frames.md!}. \textbf{Objectives:} map
\passthrough{\lstinline!Doc!}, \passthrough{\lstinline!Lib!},
\passthrough{\lstinline!Modules!}, \passthrough{\lstinline!Objects!},
\passthrough{\lstinline!Python!}, \passthrough{\lstinline!Parser!},
\passthrough{\lstinline!Include!} and \passthrough{\lstinline!Lib/test!}
responsibilities.

\textbf{Lab:} trace a public feature from documentation to Python or C
implementation and its tests. \textbf{Exit criterion:} a source map with
paths, evidence type and gaps where the snapshot text omits file bodies.

\section{M06 --- Lexer, PEG parser and
AST}\label{m06-lexer-peg-parser-and-ast}

\textbf{Prerequisite:} M05 and M01. \textbf{Reading:} Volume 04;
\path{InternalDocs/parser.md},
\path{changing_grammar.md},
\path{Doc/reference/grammar.rst}.
\textbf{Objectives:} explain PEG ordered choice and grammar generation,
inspect tokens/AST, identify syntax errors versus semantic errors.

\textbf{Lab:} build a token/AST explorer for five small programs,
including a comprehension and multiline expression. \textbf{Exit
criterion:} explain at least three AST nodes and one lexical effect; no
unsupported claim about omitted grammar bodies.

\section{M07 --- Compiler, CFG and code
objects}\label{m07-compiler-cfg-and-code-objects}

\textbf{Prerequisite:} M06 and M02. \textbf{Reading:} Volume 04;
\path{InternalDocs/compiler.md},
\path{code_objects.md},
\path{exception_handling.md}. \textbf{Objectives:}
reconstruct token-to-code stages; identify control-flow and exception
constraints; use \passthrough{\lstinline!dis!} as build-specific
evidence.

\textbf{Lab:} compare an AST and code object for a closure with
\passthrough{\lstinline!try/finally!}. \textbf{Exit criterion:} a stage
table explaining data retained at each stage and a short
semantic-equivalence argument for a hypothetical optimization.

\section{M08 --- Interpreter and optimization
architecture}\label{m08-interpreter-and-optimization-architecture}

\textbf{Prerequisite:} M07. \textbf{Reading:} Volume 03;
\path{InternalDocs/interpreter.md},
\passthrough{\lstinline!stackrefs.md!},
\passthrough{\lstinline!jit.md!}. \textbf{Objectives:} distinguish code
objects from frames; explain dispatch, specialization and optional trace
JIT at a conceptual level.

\textbf{Lab:} warm a function, inspect disassembly, measure basic
performance and record the full interpreter/build metadata. \textbf{Exit
criterion:} a technically correct discussion of why opcode output and
JIT presence are not portable language guarantees.

\section{M09 --- Reference ownership, GC, and
concurrency}\label{m09-reference-ownership-gc-and-concurrency}

\textbf{Prerequisite:} M08, M04. \textbf{Reading:} Volume 05;
\path{InternalDocs/garbage_collector.md},
\passthrough{\lstinline!qsbr.md!},
\path{Doc/c-api/refcounting.rst},
\path{Doc/howto/free-threading-python.rst}.
\textbf{Objectives:} draw ownership lifetimes, explain cycles, reason
about thread synchronization and deferred reclamation.

\textbf{Lab:} inspect a reference cycle and a synchronized shared
counter; compare allocation traces across workloads. \textbf{Exit
criterion:} an owned/borrowed reference ledger, a separate logical race
analysis, and accurate limitations of
\passthrough{\lstinline!Py\_REFCNT!}.

\section{M10 --- Native extension
implementation}\label{m10-native-extension-implementation}

\textbf{Prerequisite:} M07, M09 and M04. \textbf{Reading:} Volume 06;
\path{Doc/c-api/intro.rst},
\passthrough{\lstinline!stable.rst!},
\path{extension-modules.rst},
\path{refcounting.rst},
\path{Doc/howto/free-threading-extensions.rst}.
\textbf{Objectives:} compile a native module, check error conventions
and determine ABI/threading compatibility.

\textbf{Lab:} build the demonstration extension with a conventional
\passthrough{\lstinline!PyInit\_*!} entry point, add bad-input tests,
create ownership ledger and outline a multi-phase alternative.
\textbf{Exit criterion:} no mismatched reference ownership, no missing
error propagation and an explicit statement of unsupported
configurations.

\section{M11 --- Diagnostics, profiling and
performance}\label{m11-diagnostics-profiling-and-performance}

\textbf{Prerequisite:} M08--M10. \textbf{Reading:} Volumes 03, 05 and
07; \path{InternalDocs/profiling_binary_format.md},
\path{Doc/library/tracemalloc.rst}.
\textbf{Objectives:} distinguish Python allocations from RSS, read
tracebacks, measure variance and avoid instrument-induced false
conclusions.

\textbf{Lab:} collect a baseline profile and memory snapshot for a
constrained workload, then change one factor. \textbf{Exit criterion:}
result includes environment metadata, multiple samples, clear metrics
and no claims of guaranteed speedups.

\section{M12 --- Capstone engineering
review}\label{m12-capstone-engineering-review}

\textbf{Prerequisite:} all modules. \textbf{Objectives:} assemble an
interoperable Python application with an optional native extension,
architecture notes, threat model, reproducible tests, and a
release/rollback plan.

\textbf{Capstone specification:} implement a bounded batch processor
that reads a validated structured input file, performs a deterministic
transform, writes output safely, and reports sanitized metrics. Expose
one optional operation through a native module and fall back to pure
Python if unavailable. The capstone must run even when the optional
extension is absent.

\textbf{Exit criterion:} runnable source and tests; source-to-runtime
architecture map; ownership and synchronization review; measured
benchmark; security controls; evidence-linked release packet.

\chapter{Detailed capstone acceptance
tests}\label{detailed-capstone-acceptance-tests}

\begin{longtable}[]{@{}
  >{\raggedright\arraybackslash}p{(\columnwidth - 4\tabcolsep) * \real{0.3333}}
  >{\raggedright\arraybackslash}p{(\columnwidth - 4\tabcolsep) * \real{0.3333}}
  >{\raggedright\arraybackslash}p{(\columnwidth - 4\tabcolsep) * \real{0.3333}}@{}}
\toprule\noalign{}
\begin{minipage}[b]{\linewidth}\raggedright
Test ID
\end{minipage} & \begin{minipage}[b]{\linewidth}\raggedright
Condition
\end{minipage} & \begin{minipage}[b]{\linewidth}\raggedright
Expected result
\end{minipage} \\
\midrule\noalign{}
\endhead
\bottomrule\noalign{}
\endlastfoot
CAP-001 & Valid structured input & Correct deterministic output \\
CAP-002 & Missing required field & Explicit validation error \\
CAP-003 & Oversized input & Rejected before unbounded work \\
CAP-004 & Malformed text encoding & Clean failure with no partial
result \\
CAP-005 & Destination write failure & Original output remains valid
where supported \\
CAP-006 & Task cancelled during I/O & Resources closed and state
accounted for \\
CAP-007 & Native extension unavailable & Pure-Python fallback works \\
CAP-008 & Native method invalid operand & Python exception; no crash \\
CAP-009 & Concurrent requests & Documented state invariant preserved \\
CAP-010 & Secrets in config & No secret value appears in logs \\
CAP-011 & Profile execution & Interpreter/build metadata recorded \\
CAP-012 & CI intentionally failing test & No release artifact
published \\
\end{longtable}

\chapter{Grading and evidence
expectations}\label{grading-and-evidence-expectations}

Use a total of 100 points: \textbf{correctness 30},
\textbf{source-traced architecture 20}, \textbf{security and fault
handling 20}, \textbf{native ownership/concurrency 15},
\textbf{repeatability/performance 10}, \textbf{documentation and
operational clarity 5}. Require a passing correctness suite and no
unmitigated critical native-memory or untrusted-deserialization issues
even if the numerical total exceeds the passing threshold.

Evidence packet structure (recommended):

\begin{lstlisting}
capstone/
  src/                    # pure-Python implementation
  native/                 # optional extension
  tests/                  # unit, negative and integration tests
  docs/architecture.md    # source paths and trust boundaries
  docs/ownership.md       # native lifetime analysis
  docs/security.md        # threat model and residual risks
  benchmarks/             # methods and observed results
  requirements/           # supported runtime/packaging constraints
  README.md               # build, run, test and recovery
\end{lstlisting}

\chapter{Faculty and reviewer
guidance}\label{faculty-and-reviewer-guidance}

Distinguish \textbf{a correct observed experiment} from \textbf{a
general language rule}. Require exact citations for architectural
statements such as PEG behavior, the compiler stages, internal frame and
stack-reference concepts, free-threading guidance, and ABI tiers. If a
named source path is only present in the directory listing and not as a
body in the provided snapshot, mark it as a follow-up inspection target,
not verified source code.

Use a prerequisite graph to reschedule failed modules: a learner who
misses the ownership ledger in M09 repeats ownership-focused exercises
before attempting M10, regardless of aggregate score. A learner who
mischaracterizes \passthrough{\lstinline!pickle!} trust boundaries
repeats M04 before capstone approval.

\chapter{Evidence matrix for the eight-volume
series}\label{evidence-matrix-for-the-eight-volume-series}

\begin{longtable}[]{@{}
  >{\raggedright\arraybackslash}p{(\columnwidth - 4\tabcolsep) * \real{0.3333}}
  >{\raggedright\arraybackslash}p{(\columnwidth - 4\tabcolsep) * \real{0.3333}}
  >{\raggedright\arraybackslash}p{(\columnwidth - 4\tabcolsep) * \real{0.3333}}@{}}
\toprule\noalign{}
\begin{minipage}[b]{\linewidth}\raggedright
Volume
\end{minipage} & \begin{minipage}[b]{\linewidth}\raggedright
Curriculum coverage
\end{minipage} & \begin{minipage}[b]{\linewidth}\raggedright
Principal repository evidence
\end{minipage} \\
\midrule\noalign{}
\endhead
\bottomrule\noalign{}
\endlastfoot
01 & M01--M02 & \path{Doc/reference/},
\path{Doc/tutorial/} \\
02 & M03 & \passthrough{\lstinline!Doc/library/!},
\passthrough{\lstinline!Lib/!} \\
03 & M05, M08 & \path{InternalDocs/interpreter.md},
\passthrough{\lstinline!frames.md!}, \passthrough{\lstinline!jit.md!} \\
04 & M06--M07 & \path{InternalDocs/parser.md},
\passthrough{\lstinline!compiler.md!},
\path{code_objects.md} \\
05 & M09 & \path{InternalDocs/garbage_collector.md},
\passthrough{\lstinline!qsbr.md!} \\
06 & M10 & \passthrough{\lstinline!Doc/c-api/!},
\path{Doc/howto/free-threading-extensions.rst} \\
07 & M04, M11 &
\path{Doc/library/security_warnings.rst},
\passthrough{\lstinline!Lib/test/!} \\
08 & M01--M12 & Original prerequisite and acceptance design \\
\end{longtable}

\chapter{Source discipline and update
cadence}\label{source-discipline-and-update-cadence}

This curriculum is derived from the uploaded \textbf{3.16.0 alpha 1+dev}
text snapshot, not a final release. Each semester or internal cohort
should freeze its interpreter version, verify reference paths, refresh
removed/deprecated APIs, and maintain a change log. Compiler/JIT and
free-threading labs deserve especially strict version pinning because
they involve internals and evolving behavior.

\textbf{Source locations:} \passthrough{\lstinline!README.rst!};
\path{InternalDocs/README.md},
\passthrough{\lstinline!structure.md!},
\passthrough{\lstinline!compiler.md!},
\passthrough{\lstinline!parser.md!},
\path{interpreter.md},
\passthrough{\lstinline!jit.md!}, \passthrough{\lstinline!qsbr.md!},
\path{garbage_collector.md};
\path{Doc/reference/};
\passthrough{\lstinline!Doc/c-api/!};
\passthrough{\lstinline!Doc/howto/!};
\passthrough{\lstinline!Doc/library/!};
\passthrough{\lstinline!Lib/test/!}.

\textbf{Authorship boundary:} Module ordering, schedule, rubrics,
project architecture and tests are new teaching material. The CPython
repository does not prescribe this twelve-module learning program.

\end{document}
